\section*{Résumé}
\addcontentsline{toc}{section}{Résumé}

Au Sénégal, les enfants demeurent privés sur un ensemble de dimensions qui détériorent leur niveau de bien-être. Plus de la moitié des enfants sont touchés par la pauvreté multidimensionnelle en 2018. Parallèlement, plus de xx\% des ménages reçoivent régulièrement des transferts de migrants. L'objectif de cette étude est d'évaluer
l'impact causal des transferts
de migrants sur la pauvreté multidimensionnelle des enfants. Les deux vagues de l'Enquête harmonisée sur les conditions de vie des ménages
(EHCVM~I, 2018-2019 et EHCVM~II, 2021-2022) sont mobilisées. La pauvreté des enfants est appréhendée à travers l'approche multidimensionnelle (MODA) de l'UNICEF à partir de sept (7) dimensions réparties en 14 indicateurs. Une stratégie quasi-expérimentale basée sur le score de propension (PSM) et la double différence (DD)
est mobilisée. Trois catégories d'âge ont été construites (O-4 ans; 5-14 ans; 15-17 ans) pour capturer l'impact différencié tout au long du développement de l'enfant. Les résultats descriptifs montrent que la proportion d'enfants privés dans au moins quatre (4) dimensions est passée de 58,5\,\% en 2018-2019 à 50,1\,\% en 2021-2022.

\textbf{fournir l'impact global , par dimension significatif, selon le seuil de privation, principales implications}


 La stratégie d'identification
combine l'appariement par score de propension et la double différence
(estimateur PSM-DD de Heckman et al., 1997-1998)~: les ménages « entrants »,
sans transfert international en 2018 et bénéficiaires en 2021, sont comparés à
des ménages jamais bénéficiaires de profil comparable, ce qui contrôle
simultanément la sélection sur observables et les hétérogénéités inobservables
stables dans le temps.

L'effet moyen du traitement sur les traités est proche de zéro et non
significatif sur l'incidence N-MODA (ATT~=~-0,011~; $p = 0{,}676$)~: sur un
horizon de trois ans, la privation des enfants des ménages devenus
bénéficiaires a reculé au même rythme que celle des enfants de ménages
comparables jamais bénéficiaires, conclusion confirmée par la double différence
brute. Ce résultat nul reflète le délai de matérialisation des transferts sur
des privations structurelles, la compensation entre l'apport monétaire et les
coûts de réorganisation du ménage consécutifs au départ du migrant, et les
contraintes d'offre de services de base que le seul pouvoir d'achat ne lève
pas.

La décomposition par dimension en apporte la preuve directe~: la privation
éducative diminue significativement chez les entrants (-3,9\,pp~;
$p = 0{,}005$), tandis que la protection de l'enfant se dégrade en sens inverse
(5,1\,pp~; $p = 0{,}003$), deux canaux opposés qui se neutralisent dans
l'agrégat. La définition alternative fondée sur une exposition durable (ménages
bénéficiaires aux deux vagues) révèle par ailleurs une dynamique relative
défavorable, quoique non significative (0,039~; $p = 0{,}108$), suggérant que les effets
se manifestent au-delà de l'horizon de trois ans du design principal.

À court terme, commencer à recevoir des transferts ne modifie donc pas la
trajectoire de privations multidimensionnelles des enfants sénégalais. Ces
résultats plaident pour des politiques combinant la réduction des coûts de
transaction des envois, l'accompagnement des ménages dans la durée et le
renforcement de la complémentarité entre transferts privés et services publics
de base.

\bigskip

\noindent \textbf{Mots-clés} : pauvreté multidimensionnelle des enfants, N-MODA,
transferts de migrants, Sénégal, EHCVM, appariement par score de propension (PSM),
double différence (DD).
